\section{Model}
\label{sec:model}

\subsection{Circuits and tensor networks}

\begin{definition}[Circuit or tensor network]
A circuit or tensor network \(C\) consists of:
\begin{enumerate}[leftmargin=1.2cm]
\item a finite set \(V(C)\) of tensors, including input states, gates, channels represented by tensors, effects, measurements, discards, and optional postselection tensors;
\item a finite set \(I(C)\) of indices, each with dimension \(d_i\), where internal indices are contracted and boundary indices are fixed, summed, or measured;
\item a designated output event \(x\), such as a computational-basis outcome on a subset of output wires.
\end{enumerate}
The scalar \(p_C(x)\) denotes the probability of observing \(x\), computed by the Born rule or equivalently by contracting the doubled tensor network for the event \(x\).
\end{definition}

The formalism includes ordinary quantum circuits, open circuits with discards, tensor-network representations of effects, and postselected networks provided the target scalar is well-defined. For algorithmic statements, we assume local dimension is bounded by a constant, typically \(d_i=2\), unless explicitly stated.

\begin{definition}[Compositional quantum language circuit]
A compositional quantum language circuit family is a family \(\{C_\alpha\}_{\alpha\in A}\) indexed by grammar, discourse, or document composition data \(\alpha\). The index \(\alpha\) determines a composition diagram \(D_\alpha\), and \(C_\alpha\) is obtained by assigning local tensors or circuit fragments to the generators and wires of \(D_\alpha\), followed by a fixed translation to a quantum circuit or tensor network.
\end{definition}

This definition is deliberately broad. It includes pregroups and string diagrams, lambeq-generated ansatz circuits, QDisCoCirc-style discourse circuits, and document-level tensor networks obtained by composing sentence or paragraph modules.

\subsection{Contraction graph}

\begin{definition}[Contraction graph]
The contraction graph \(G_C=(V_C,E_C)\) of \(C\) has one vertex for each tensor in \(V(C)\). For each contracted index shared by two tensors, \(G_C\) has an edge between the corresponding vertices. If an index is shared by more than two tensors, one may either use a hyperedge or introduce an equality/copy tensor and work with an ordinary graph. Boundary indices are represented as dangling legs or as edges to fixed boundary tensors.
\end{definition}

For quantum circuits, \(G_C\) is not necessarily the circuit interaction graph on qubits. It is the graph relevant to tensor contraction of the scalar \(p_C(x)\). The doubled network for probabilities may have a different graph from the amplitude network; the theorem statements apply to whichever contraction graph is used by the simulator.

\subsection{Tree decompositions and width}

\begin{definition}[Tree decomposition]
A tree decomposition of \(G_C\) is a pair \((T,\beta)\), where \(T\) is a tree and \(\beta:V(T)\to 2^{V_C}\) assigns a bag of graph vertices to each tree node, such that:
\begin{enumerate}[leftmargin=1.2cm]
\item for every \(v\in V_C\), the set \(\{b\in V(T):v\in\beta(b)\}\) is nonempty and connected in \(T\);
\item for every edge \(\{u,v\}\in E_C\), there exists \(b\in V(T)\) with \(u,v\in\beta(b)\).
\end{enumerate}
The vertex width is \(w_{\mathrm{v}}=\max_b |\beta(b)|-1\).
\end{definition}

Because tensor contraction cost depends on index dimensions rather than only the number of graph vertices in a bag, we also use an index-separator width.

\begin{definition}[Separator width]
Root \(T\) at an arbitrary bag \(r\). For an oriented edge \(e=(b\to c)\), let \(S_e\) be the set of tensor indices with one endpoint in the subtree of \(c\) and the other endpoint outside it, represented in the contracted network. The separator width is
\[
w=\max_{e\in E(T)} \sum_{i\in S_e}\log_2 d_i .
\]
For bounded local dimension and bounded tensor degree this is equivalent, up to constants, to the usual treewidth-type contraction width.
\end{definition}

\subsection{Free and resource tensors}

\begin{definition}[Stabilizer/free tensor]
A tensor is stabilizer/free if it belongs to the tensor class generated by computational-basis states, stabilizer states, Clifford unitaries, Pauli measurements, Clifford effects, tensor product, composition, wire bending allowed by the chosen graphical calculus, and partial trace/discard operations that preserve efficient stabilizer simulation.
\end{definition}

This definition is intentionally operational. A free subnetwork is one for which the relevant contracted messages can be represented and updated in polynomial time by stabilizer-tableau, affine symplectic, or equivalent stabilizer data structures, up to the explicit active-boundary factor \(2^{O(\weff)}\) when a dense active table is required.

\begin{definition}[Resource tensor]
A resource tensor is a tensor not included in the free stabilizer class after maximal local Clifford simplification, Pauli propagation, stabilizer-state absorption, and contraction of adjacent free components. The resource set is denoted \(\R(C)\subseteq V_C\).
\end{definition}

The phrase ``maximal local Clifford simplification'' is not a canonical mathematical operation unless a normalization procedure is specified. Here it means any sound preprocessing that reduces free components without increasing the chosen local magic costs. A fully specified implementation fixes such a procedure or quantifies over all valid reductions.

\subsection{Output probabilities and doubled networks}

For a unitary circuit with input \(\rho_{\mathrm{in}}\), measurement POVM element \(M_x\), and implemented channel \(\Phi_C\), the target probability is
\[
p_C(x)=\Tr[M_x\Phi_C(\rho_{\mathrm{in}})].
\]
In tensor-network notation this may be represented either as a channel network with open input/output legs paired by trace tensors, or as a doubled ket-bra network. The choice affects the contraction graph and hence the width parameter. A simulation theorem must therefore specify the graph used for the scalar being estimated. Throughout this paper, \(G_C\) denotes the graph after this choice has been made.

For postselected circuits, let \(s\) denote the postselection event and \(x\) the reported outcome. One can estimate the joint scalar \(p_C(x,s)\) and, when \(p_C(s)\) is not too small and is also estimable, form the conditional probability \(p_C(x\mid s)\). The main theorem is stated for a single scalar \(p_C(x)\). Extending the bound to ratios requires standard error propagation and a lower bound on the denominator.

\subsection{Uniformity and input model}

The algorithmic statements assume that the family \(\{C_n\}\) is classically describable in time \(\poly(n)\), local tensor dimensions are bounded or explicitly included in the width, and local free-tensor operations admit polynomial-time stabilizer updates. Tree decompositions may be supplied as input or computed by a heuristic. If a decomposition is computed, its cost is added to the runtime. Since exact treewidth is NP-hard in general, the theorem is an algorithmic guarantee conditional on a decomposition of the stated width.

\subsection{Equivalence transformations}

Several transformations preserve the target scalar while changing the apparent resource distribution:
\begin{itemize}[leftmargin=1.2cm]
\item Clifford gates can be propagated through adjacent Clifford regions and absorbed into stabilizer states or measurements.
\item Adjacent inverse gates or cups/caps in a compact-closed diagram may cancel.
\item Free tensors inside a subtree can be contracted into a boundary stabilizer message.
\item A non-Clifford unitary can be represented as injection of a magic state followed by a Clifford gadget, moving the resource from a gate tensor to a state tensor.
\end{itemize}
The separator-local burden is minimized over such sound transformations only when they are included in the simulator. Otherwise the parameter would describe an unavailable contraction strategy.
